\section{Síntesis y ampliación}

\begin{enumerate}
    \item Offline RL optimiza la política sobre datos estáticos y necesita equilibrar la imitación y la exploración.
    \item La cobertura del conjunto de datos y la calidad de reward determinan directamente la robustez del método.
    \item IQL restringe implícitamente la búsqueda mediante expectile V, mientras que CQL impone conservatism mediante explicit penalty.
    \item La evaluación multinivel, el monitoreo mediante logging y el marco de gobernanza son clave para un despliegue exitoso.
    \item Monitorear el shift de distribución, Q variance y action entropy permite controlar el riesgo tras el despliegue.
\end{enumerate}

\subsection{Tabla de síntesis}

\begin{table}[H]
\centering
\begin{tabular}{p{0.2\textwidth} p{0.24\textwidth} p{0.24\textwidth} p{0.2\textwidth}}
\toprule
Dimensión & Problema central & Medida de respuesta & Herramienta típica \\
\midrule
Conjunto de datos & Cobertura, calidad de la recompensa & Filtrar recompensas bajas, ponderar imitation & D4RL, dataset diagnostics \\
Algoritmo & Un desplazamiento grande provoca una sobreestimación de Q & IQL expectile, CQL penalty & IQL codebase, CQL repo \\
Evaluación & benchmark score vs deployment performance & Evaluación multinivel, human audit & D4RL + replay evaluation + human review \\
Despliegue & Monitoreo de drift, reward shift & logging/alert + replay & dashboards, entropy targets \\
Gobernanza & traceability + sign-off & records, policy reviews & notebooks, compliance logs \\
Tooling & pipeline automation & data-versioning, policy registry & W\&B, Feathr, Kubeflow \\
Research & offline→online handoff safety & counterfactual auditing & world models, meta RL supervisors \\
\bottomrule
\end{tabular}
\caption{Tabla de control de ingeniería para proyectos de Offline RL}
\end{table}

\subsection{Acciones adicionales}

\begin{enumerate}
    \item Usar primero la tabla de dataset diagnostics para identificar trayectorias buenas o malas, y luego decidir si aplicar filter, IQL o CQL.
    \item Construir un replay simulator antes del despliegue, y reproducir escenarios de alto riesgo mediante replay.
    \item Establecer un monitor dashboard (action KL, reward gap, Q variance) y definir claramente el alert threshold para el devops team.
    \item Registrar la audit info de cada actualización de policy (data version, reward log, human sign-off) en el registro de gobernanza.
\end{enumerate}

\subsection{Lecturas adicionales}
\begin{itemize}
    \item Levine et al., \textit{Offline Reinforcement Learning: Tutorial, Review, and Perspectives on Open Problems} (2020).
    \item Kostrikov et al., \textit{Offline Reinforcement Learning with Implicit Q-Learning} (IQL, 2022).
    \item Kumar et al., \textit{Conservative Q-Learning for Offline Reinforcement Learning} (CQL, 2020).
    \item Laskin et al., \textit{Offline Reinforcement Learning with Dense Rewards for Realistic Tasks} (2023).
    \item cs224r course notes repository for practical code samples: \url{https://cs224r.stanford.edu/spring_2025/lecture/07.html}
\end{itemize}

\end{document}
